\section{Marco Teórico}
\subsection{¿Qué es un Amplificador de Audio?}
Un amplificador de audio es un circuito electrónico cuya función principal es incrementar la amplitud de una señal de entrada variable en el tiempo (normalmente proveniente de fuentes como micrófonos, reproductores de audio o instrumentos) para entregar una potencia mayor a una cierta carga, tal como un altavoz o parlante.

\subsection{Clasificación de Amplificadores y Etapa de Potencia}

Los amplificadores se clasifican según su forma de amplificar la señal de entrada, separandose en clase A, B, AB, D, entre otros. El amplificador TDA2030 es de clase AB, el cual combina aspectos de los amplificadores de clase A y B.

Al igual que en el calse B, se utiliza un par de transistores push-pull, pero se elimina el curce por cero característico del amplificador de clase B. Este tipo de amplificador mejora el rendimiento típico de un clase A (al rededor del 25\%) a un 78\% aproximadamente, manteniendo la linealidad.
